\documentclass[11pt,a4paper]{article}
\usepackage[utf8]{inputenc}
\usepackage[top=0.9in, bottom=0.8in, left=1in, right=1in]{geometry}
\usepackage{hyperref}

\setlength{\parskip}{0.8em}
\setlength{\parindent}{0pt}

\begin{document}

\begin{center}
    {\LARGE \textbf{Liuxi Mei}} \\
    \vspace{2mm}
    Linköping, Sweden $\cdot$ liume102@student.liu.se $\cdot$ +46-0704898962
    \vspace{4mm}
\end{center}

\begin{flushright}
May 18, 2026
\end{flushright}
 
Dear Professor Rosenquist Brandell,

I am writing to apply for the PhD Position in Data-Driven Precision Medicine and Diagnostics (DDLS Program) at Karolinska Institutet, within the Clinical Genetics Research Group at the Department of Molecular Medicine and Surgery. I am currently completing a second master's degree in Statistics and Machine Learning at Linköping University (expected graduation September 2026). I already hold a master's degree in Computer Science (Central South University for Nationalities, 2015).

\textbf{Motivation for the position} What makes this PhD particularly exciting to me is that it is centered on a real clinical puzzle rather than a benchmark problem. In CLL, two patients can look similar based on known drivers and still follow very different trajectories, and even with targeted therapies such as BTK and BCL2 inhibitors, resistance emerges in ways that are not always obvious from standard genetic markers. The question I want to work on is how to use the full picture of patient molecular data to explain that divergence and make it actionable.

What appeals to me is that the project is not only about building a model that predicts something, but about integrating genomics, epigenomics, transcriptomics and proteomics with clinical information and prior biological knowledge so that the output becomes a short list of plausible subtypes, biomarkers, and mechanisms that can be followed up. The knowledge graph aspect is a big part of what draws me in, because it provides a bridge between high-dimensional signals and the pathways and regulatory relationships that clinicians and biologists can reason about.

This also connects naturally to my master’s thesis work. I used graph neural networks together with curated interaction networks such as OmniPath as a structured prior, precisely to link high-dimensional gene expression patterns back to known biological relationships. That experience made me appreciate how much difference the choice of prior knowledge and the evaluation design can make when the goal is interpretation rather than only prediction.

Finally, the context matters to me. Working at BioClinicum in a team that spans bioinformatics, molecular biology and clinical science, with access to NAISS and the DDLS Research School environment, feels like the right place to do data-driven work that stays close to translational questions.

\textbf{Suitability for research topic:} Your call for strong computational skills, omics analysis experience, and an interest in AI and network-based methods fits my background well. I spent five years in industry building and maintaining data pipelines and backend systems. That is where I learned to write Python and other languages that other people can trust, and to keep work reproducible and clearly labeled.

On the data side, I have worked with high-dimensional gene expression data and curated biological network resources. While I have not previously worked with in-house patient cohorts, I am comfortable with the kind of careful data work this project will require early on. This includes consistent preprocessing across modalities, quality control, and baseline modeling before moving to more complex integration.

On the modeling side, my thesis experience with graph neural networks and curated interaction networks such as OmniPath gives me a practical starting point for using biological prior knowledge as model structure. In industry, I also did data mining, analysis and prediction on time series and map-based data.

\textbf{Professional and academic background} I spent more than five years in industry building data-intensive systems. That experience taught me how to work with messy inputs, inconsistent schemas, and systems that need to be reliable and maintainable over time. In my current master's program, I have focused on statistical learning and machine learning, and I have been moving toward biomedical applications where model design and evaluation need to be careful and honest.

\textbf{Master’s thesis research} In my thesis work, I used graph neural networks together with curated biological interaction networks (0mnipath) to combine network context with gene expression signals. The goal was to learn representations that are predictive while staying connected to known biology. The project gave me practical experience with biological data processing, model evaluation, and communicating results in a way that is useful to people outside of machine learning. It also shaped how I would approach this PhD project: start from a solid data and evaluation pipeline, build strong baselines, and then iterate toward more integrative models that improve interpretability and robustness.

Thank you for considering my application. I would be happy to discuss how my engineering experience and current research direction could support your group’s DDLS precision medicine work in CLL.

Sincerely,

Liuxi Mei

\end{document}
